\begin{pastexam}[label={px:ncku-113-cce-13}]{成大電通所 113 學年度 簡答題第 4 題}
（簡答）

(6 pts) Please show the data structures to record the following graph in the form of adjacency
matrix and the inversed adjacency list.

\begin{center}
\begin{tikzpicture}[>={Latex[length=5pt]},thick,v/.style={draw,circle,minimum size=0.8cm,inner sep=0pt,font=\sffamily\large}]
  \node[v] (0) at (0,0) {0}; \node[v] (1) at (1.5,0) {1}; \node[v] (2) at (3.1,0) {2};
  \draw[->] ([yshift=0.22cm]1.west)--([yshift=0.22cm]0.east);
  \draw[->] ([yshift=-0.22cm]0.east)--([yshift=-0.22cm]1.west);
  \draw[->] (1)--(2);
\end{tikzpicture}
\end{center}
\end{pastexam}

\begin{pastsoln}
圖有三條有向邊：$0 \to 1$、$1 \to 0$、$1 \to 2$。

\textbf{相鄰矩陣}：$A[i][j] = 1$ 表示有邊 $i \to j$，否則是 $0$（資結書第~21.2~節）。 第 $i$ 列
是 $i$ 的出邊，第 $j$ 行是 $j$ 的入邊。
\begin{center}
\begin{tabular}{c|ccc}
 & 0 & 1 & 2 \\ \hline
0 & 0 & 1 & 0 \\
1 & 1 & 0 & 1 \\
2 & 0 & 0 & 0 \\
\end{tabular}
\end{center}
有向圖的矩陣不一定對稱：$A[1][2] = 1$ 但 $A[2][1] = 0$。

\textbf{反鄰接串列}（inverse adjacency list）：每個頂點 $v$ 一條串列，存所有\textbf{指向} $v$
的頂點 $u$（有邊 $u \to v$），也就是相鄰矩陣第 $v$ 行裡的 $1$。
\begin{center}
\begin{tikzpicture}[>={Latex[length=4pt]},
  hd/.style={draw,minimum width=0.6cm,minimum height=0.6cm,inner sep=0pt},
  nd/.style={draw,minimum width=0.7cm,minimum height=0.6cm,inner sep=0pt},
  lk/.style={draw,minimum width=0.5cm,minimum height=0.6cm,inner sep=0pt}]
  \foreach \v/\u/\y in {0/1/0, 1/0/-0.6, 2/1/-1.2} {
    \node at (-0.75,\y) {$\v$};
    \node[hd] (h\v) at (0,\y) {};
    \node[nd] (n\v) at (2,\y) {\u};
    \node[lk,anchor=west] (l\v) at (n\v.east) {};
    \draw[->] (h\v.center)--(n\v.west);
    \fill (h\v.center) circle (1.2pt);
    \draw (l\v.south west)--(l\v.north east);
  }
\end{tikzpicture}
\end{center}
即 $0\colon 1$；$1\colon 0$；$2\colon 1$（斜線表示 \NIL）。 用反鄰接串列，找一個頂點的所有入邊
只要走它自己那一條串列，入度就是串列長度；一般的鄰接串列則要掃過整個結構才找得到入邊。
\end{pastsoln}
